\section{Trabajos Relacionados}
\label{sec:related_work}

La inconsistencia y carencia de normalización en conjuntos de imágenes de resonancia magnética obstaculizan tanto el diagnóstico automatizado como la precisión en la detección de patologías neurológicas. En una investigación previa, se propuso valorar el desempeño comparativo de diversas arquitecturas preentrenadas mediante la estrategia de aprendizaje por transferencia para categorizar cuatro tipos distintos de masas cerebrales. El aporte fundamental consistió en analizar de forma sistemática arquitecturas profundas tales como Xception, ResNet50, VGG16 y DenseNet121. La aproximación metodológica utilizó un conjunto de datos público disponible en Kaggle que contenía 7023 imágenes MRI dedicadas al entrenamiento y optimización de los parámetros de la red. El modelo Xception logró una precisión del $98.73\%$ como resultado principal. Sin embargo, se identificó una limitación importante: menor sensibilidad para las categorías glioma y meningioma, \cite{Disci2025}.

La degradación de resolución espacial y la presencia de artefactos técnicos en resonancias magnéticas reducen la eficiencia en la detección precoz de neoplasias. Un estudio previo se orientó a mejorar los procesos de categorización mediante el refinamiento secuencial específico de la familia de redes EfficientNet. La contribución radió en la explotación sistemática de los parámetros de escalamiento de EfficientNetB0-B4 orientada a capturar descriptores sutiles de los tejidos. La metodología incorporó el procesamiento de un dataset público proveniente de Figshare junto con un flujo de trabajo de entrenamiento que incluía optimizadores avanzados para prevenir sobreajuste. Se consiguió un $99.06\%$ de accuracy, $98.73\%$ en precisión, $99.13\%$ en sensibilidad y un F1-score de $98.79\%$. Como deficiencia notable, se carece de un protocolo de validación externa bajo condiciones clínicas auténticas \cite{BabuVimala2023}.

El consumo computacional elevado y la pérdida de información de escala global disminuyen la capacidad de las redes neuronales independientes en problemas de categorización multiclase de tumores. Una investigación desarrolló un sistema híbrido de clasificación que sintetiza los beneficios de múltiples extractores automáticos de características. El aporte radicó en el diseño de un modelo ensamblado que integra componentes de ResNet50V2, MobileNetV2 y DenseNet121. El procedimiento metodológico fusionó los vectores característicos derivados de 7023 imágenes de un repositorio público antes del módulo de clasificación terminal. El modelo resultante logró un $98.75\%$ de accuracy, con métricas de precisión del $98.76\%$ y F1-score del $98.75\%$. Como debilidad identificada, se utilizaron exclusivamente datos públicos preprocesados, sin considerar el impacto que el ruido inherente al ambiente clínico y los valores atípicos ejercen sobre sistemas complejos \cite{DwiPrayogo2025}.

Las variaciones morfológicas complejas presentes en imágenes de resonancia magnética limitan la discriminación de redes convolucionales tradicionales para diferenciar entre múltiples estadios de masas tumorales. Una investigación previa optimizó la categorización multiclase mediante el ajuste fino secuencial de ResNet50 enriquecida con capas totalmente conectadas suplementarias. El aporte esencial consistió en el diseño de un módulo de clasificación personalizado que maximiza la preservación de atributos jerárquicos preentrenados. El procedimiento requirió la consolidación y procesamiento de tres bases de datos públicas distintas: SARTAJ, Br35H y Figshare. Como resultado, la arquitectura modificada alcanzó una exactitud sobresaliente del $99.31\%$, con mejoras paralelas en F1-score, precisión, sensibilidad y el área bajo la curva (AUC). La brecha metodológica identificada radica en que no se realizó contraste del rendimiento mediante validación externa clínica \cite{Mitra2025}.

La limitación intrínseca de las operaciones convolutivas de alcance local para modelar relaciones espaciales de largo rango afecta la precisión en la delineación de bordes tumorales difusos. Una investigación exploró la viabilidad de mecanismos de auto-atención global mediante la aplicación del modelo preentrenado Vision Transformer ViT-B/16. La contribución principal residió en la explotación del procesamiento parcelado de imágenes (patch embedding) para capturar vínculos contextuales sofisticados en cortes encefálicos. El procedimiento llevó a cabo un refinamiento integral sobre un conjunto de datos público de resonancia magnética estructurado en cuatro categorías patológicas. El experimento reportó una exactitud máxima del $99.24\%$ en el conjunto de validación para ViT-B/16. La brecha metodológica identificada consiste en que el estudio no especifica el empleo de una cohorte externa de validación \cite{Panigrahi2025}.

El costo operacional significativo y la demanda masiva de datos que caracterizan a los Transformers convencionales limitan su viabilidad práctica en entornos de diagnóstico médico automatizado. Una investigación llevó a cabo una comparación directa entre las capacidades de procesamiento del Vision Transformer estándar y la arquitectura Swin Transformer con ventanas desplazables. La aportación central se enfocó en demostrar la eficiencia del Swin Transformer en el procesamiento de información a múltiples resoluciones espaciales con complejidad computacional lineal. El procedimiento implementó el entrenamiento estratificado de ambas aproximaciones utilizando un repositorio público indexado en Kaggle que contenía un total de 3000 imágenes MRI. Los resultados estadísticos corroboraron la superioridad del modelo Swin al alcanzar un $98.71\%$ de exactitud frente al $95.30\%$ logrado por la arquitectura ViT. Como limitación, el estudio se restringió a un dataset de escala moderada \cite{Jain2025}.

La pérdida de información de alta frecuencia causada por el escalamiento global de los Transformers afecta la precisión en la delimitación de anomalías de pequeña magnitud milimétrica. Un trabajo previo integró las capacidades locales de las redes convolucionales con el modelado de dependencias globales característico de módulos Transformer dentro de una estructura unificada. El aporte consistió en el desarrollo del modelo híbrido TECNN, que funciona como equivalente a una arquitectura CvT (Convolutional Vision Transformer). El procedimiento validó de forma rigurosa el pipeline computacional usando dos plataformas públicas de manera independiente: BraTS 2018 y Figshare. El sistema logró un $96.75\%$ de exactitud en el dataset desafiante de BraTS 2018 y un $99.10\%$ en Figshare. Su limitación radica en que el marco no emplea la arquitectura CvT nominal de referencia, sino una aproximación híbrida propia \cite{Aloraini2023}.

La degradación de gradientes en redes profundas obstaculiza la extracción simultánea de patrones locales de textura y conceptos contextuales abstractos. Un estudio diseñó un marco de aprendizaje simétrico de dos ramas paralelas que integra módulos Transformer con redes residuales optimizadas. La aportación se basó en el acoplamiento de un módulo de auto-atención (Self-Attention) junto a una arquitectura iResNet adaptada para estabilizar la propagación de información. El procedimiento requirió el análisis de 3064 secciones de resonancias magnéticas del repositorio Figshare, distribuidas en cuatro categorías diagnósticas. Los ensayos de validación demostraron que el extractor basado en iResNet alcanzó un $99.30\%$ de exactitud, mejorando la capacidad de generalización en un $20\%$ de muestras no observadas durante el ajuste. La debilidad identificada consiste en la carencia de validación frente a una cohorte clínica externa auténtica \cite{Tabatabaei2023}.

La separación metodológica entre sistemas de segmentación y clasificación obliga a las instituciones médicas a desplegar pipelines computacionales redundantes y onerosos. Una investigación unificó ambas labores dentro de un único marco de aprendizaje optimizado para datos médicos. La aportación principal radicó en el diseño de un Adaptive Masked Transformer capaz de resolver tareas múltiples de forma simultánea. El procedimiento empleó 485 casos clínicos privados junto con un conjunto de datos independiente para la validación externa del sistema. El modelo alcanzó un Dice de $0.91$ y una exactitud de clasificación del $93.4\%$. La limitación del enfoque consiste en que la precisión final de la categorización quedó considerablemente por debajo de los modelos especializados dedicados a etiquetado \cite{Lv2025}.

La prevalencia de ruido de alta frecuencia y la variabilidad de contraste entre escáneres de distintos fabricantes deterioran el rendimiento de las redes U-Net convencionales. Una investigación estructuró un flujo de trabajo progresivo que integrara filtrado de ruido, segmentación y posterior categorización jerárquica de gliomas. El aporte residió en un proceso de trabajo integrado que conecta una red U-Net con clasificadores derivados de VGG y GoogLeNet. El procedimiento validó el pipeline sobre el conjunto BraTS2019, ejecutando tareas de segmentación y categorización en etapas secuenciales. Los resultados alcanzaron un Dice de $0.82 / 0.91 / 0.72$ en las subtareas de segmentación y un $97.44\%$ de exactitud en la categorización. La limitación radica en que el modelo requiere mejoras arquitectónicas fundamentales por carencia de originalidad en los bloques de integración \cite{Dang2022}.

Al revisar la literatura científica contemporánea, se identifican dos brechas metodológicas de relevancia fundamental. Primero, la mayor parte de las investigaciones previas orientadas a datasets públicos de Kaggle o Figshare \cite{Disci2025, BabuVimala2023, DwiPrayogo2025, Mitra2025, Panigrahi2025, Jain2025} incurren en el error técnico de aplicar algoritmos de expansión sintética de datos de forma global sobre la totalidad del conjunto de imágenes previo a la partición. Esta práctica causa un fenómeno grave de fuga de datos (Data Leakage), donde variantes modificadas de una imagen original terminan distribuidas simultáneamente en los conjuntos de entrenamiento, validación y prueba, inflando artificialmente las métricas de desempeño y reduciendo la capacidad genuina de generalización.

Segundo, la generalidad de los trabajos previos abordan el problema como categorización multiclase (glioma, meningioma, pituitario, sin tumor). Aunque esta granularidad diagnóstica posee valor, introduce una complejidad añadida que dificulta la evaluación controlada de pipelines de preprocesamiento y particionamiento. Por este motivo, el presente trabajo adopta un enfoque de categorización binaria (sí/no tumor), para lo cual se seleccionó un subconjunto aleatorio del 20\% de los datos originales y se reestructuró explícitamente en dos categorías: sí (con tumor, integrando gliomas, meningiomas y demás variantes) y no (cerebro sano, sin tumor). Esta simplificación permite aislar y valorar con mayor rigor el impacto de las decisiones metodológicas sobre el preprocesamiento de datos, minimizando las variables confusoras asociadas a la granularidad multiclase.

Asimismo, se evidencia una carencia de protocolos sistematizados para la detección y corrección de anomalías visuales u outliers previos al entrenamiento. El presente trabajo aborda de manera directa estas carencias implementando un pipeline de automatización rigurosa (Data Automation Split) donde la partición en tres grupos independientes (Train, Validation y Test) se ejecuta de forma prioritaria y absoluta antes de cualquier técnica de expansión sintética de datos. De igual forma, se integra el uso de algoritmos avanzados de detección de anomalías (Outliers) para garantizar la pureza del conjunto de datos y asegurar la robustez clínica del modelo propuesto.

A continuación, se presenta la Tabla \ref{tab:related_work_summary}, la cual sintetiza los aportes, resultados y deficiencias metodológicas de las investigaciones analizadas en esta sección.

\begin{table*}[t]
  \caption{Tabla comparativa de la literatura científica analizada para la detección y segmentación de tumores cerebrales.}
  \label{tab:related_work_summary}
  \centering
  \small
  \setlength{\tabcolsep}{4pt}
  \renewcommand{\arraystretch}{0.9}
  \begin{tabular}{>{\raggedright\arraybackslash}p{0.13\textwidth} >{\raggedright\arraybackslash}p{0.26\textwidth} >{\raggedright\arraybackslash}p{0.22\textwidth} >{\raggedright\arraybackslash}p{0.26\textwidth}}
    \hline
    Autor y Año & Modelo / Aporte Principal & Resultados Clave & Deficiencia o Brecha Identificada \\
    \hline
    \cite{Disci2025} & Transfer Learning (Xception, ResNet50, VGG16) & Accuracy: $98.73\%$, F1: $95.29\%$ & Menor recall en glioma; explicabilidad y diversidad de datos limitadas. \\
    \cite{BabuVimala2023} & EfficientNetB0-B4 fine-tuned en CE-MRI & Accuracy: $99.06\%$, F1: $98.79\%$ & Ausencia total de validación externa en entornos clínicos reales. \\
    \cite{DwiPrayogo2025} & Ensamble híbrido (ResNet50V2 + MobileNetV2 + DenseNet121) & Accuracy: $98.75\%$, F1: $98.75\%$ & Dependencia de datasets públicos compuestos; sensible a ruido y outliers. \\
    \cite{Mitra2025} & Fine-tuning de ResNet50 + Capas adicionales & Accuracy: $99.31\%$, Alto AUC & No reporta validación externa con muestras de origen clínico real. \\
    \cite{Panigrahi2025} & Vision Transformer (ViT-B/16) fine-tuned & Accuracy: $99.24\%$ & Falta de descripción de cohortes externas y validación multicéntrica. \\
    \cite{Jain2025} & Comparación ViT vs. Swin Transformer & Swin Acc: $98.71\%$, ViT Acc: $95.30\%$ & Dataset de escala moderada y ausencia de esquemas de validación externa. \\
    \cite{Aloraini2023} & TECNN híbrido (Convolución + Transformer) & Acc BraTS: $96.75\%$, Figshare: $99.10\%$ & No emplea CvT nominal; es un desarrollo híbrido propio ad-hoc. \\
    \cite{Tabatabaei2023} & Dos ramas (iResNet + Self-Attention Module) & Accuracy: $99.30\%$, Gen. $+20\%$ & Evaluado solo en bases públicas; sin reporte de cohorte clínica externa. \\
    \cite{Lv2025} & Multitarea con Adaptive Masked Transformer & IoU: $0.921$, Clasificación Acc: $93.4\%$ & La precisión de clasificación queda por debajo de los modelos monopropósito. \\
    \cite{Dang2022} & Pipeline integrado UNet + VGG/GoogLeNet & Segmentación Dice: $0.82$; Clasificación Acc: $97.44\%$ & Falta de originalidad en la arquitectura base; requiere optimización estructural. \\
    \hline
  \end{tabular}
\end{table*}
